\section{Introduction}
\label{sec:introduction}

Peru's labor market informality rate exceeds 70\%, one of the highest in Latin America. The COVID-19 pandemic, which triggered strict lockdowns and mobility restrictions beginning in March 2020, disrupted formal employment relationships across all sectors---but its impact was not uniform. Workers in contact-intensive occupations, where physical presence is essential and remote work is infeasible, faced the most severe labor disruptions. This paper examines whether these differentially exposed workers experienced a persistent increase in informality or whether they recovered to pre-pandemic formality levels by 2024.

We exploit differential sectoral exposure to COVID-19 mobility restrictions using the \citet{dingel2020} teleworkability index as a source of identifying variation. The index, originally constructed from O*NET task descriptions for U.S. occupations, measures the feasibility of performing each occupation remotely. We map this index to Peruvian occupation codes (ISCO-08) via a crosswalk, creating a continuous, pre-determined treatment intensity variable.

Using the ENAHO rotating panel from 2020 to 2024 (Module 500, Employment), we estimate a difference-in-differences specification interacting year dummies with the continuous teleworkability score. Our main specification yields negative interaction coefficients of $-0.19$ to $-0.22$ ($p < 0.05$ in 2022 and 2024), indicating that a one-unit increase in teleworkability is associated with a 19--22 pp smaller informality increase. The binary specification (contact-intensive vs.\ teleworkable) yields a differential of approximately 9 pp, though this is not statistically significant at the 5\% level with ISCO-clustered standard errors ($p = 0.153$ in 2021). The differential persisted unchanged through 2024 (Wald test $p = 0.937$), indicating that the population-level informality gap between occupational groups has not closed over four years.

We emphasize three important caveats. First, our panel begins in 2020---the year of the shock---so we cannot test parallel pre-trends. The identifying assumption of common trends rests on plausibility, not statistical evidence. Second, ENAHO's rotating panel design limits within-individual tracking to approximately two consecutive years. Third, attrition from the panel is not differential by treatment status ($p = 0.163$), but the overall attrition rate is high (78.7\% by 2021).

These findings contribute to the literature on labor market scarring \citep{arulampalam2001} and informality dynamics in developing countries \citep{bosch2025}. The paper proceeds as follows: Section~\ref{sec:literature} reviews the literature, Section~\ref{sec:data} describes the data, Section~\ref{sec:strategy} presents the identification strategy, Section~\ref{sec:results} reports results, Section~\ref{sec:robustness} discusses robustness, and Section~\ref{sec:conclusion} concludes.
